\chapter{Discussion}
\label{ch:discussion}

\section{Interpretation of Results}
\label{sec:interpretation}

\subsection{Entity Classification}

The final V4 checkpoint reaches corrected class-aware entity micro F1 of
$0.792$ and macro F1 of $0.772$. KEY F1 is $0.720$, and VALUE F1 is $0.824$.
The calculation penalises confusion between the two entity classes: 803 KEY
targets are predicted as VALUE, and 584 VALUE targets are predicted as KEY.
The reported micro and macro aggregates include KEY and VALUE and exclude the
\texttt{O} class.
The model constrains its textual predictions to tokens in the extracted input
and therefore cannot emit text that is absent from that input, although it can
still assign an incorrect label to an observed token.

The token-level entity F1 values are numerically higher than the Regular
text+location pair F1 of $0.345$. The measures are not
directly comparable, but the higher token-level entity scores show that correct
entity recognition alone does not ensure correct relations. Together with the
observed wrong-link examples, this suggests that relation linking remains a
major internal limitation of the current pipeline.

\subsection{Token-Level Linking Failure}

V1 uses a dense token-level candidate formulation, producing a substantially
larger candidate space than the span-level formulation. The failure is
consistent with token-level class imbalance and a mismatch between token-level
candidates and span-level relations.
The illustrative ratio of approximately $500{:}1$ shows the severity of the
token-level candidate imbalance and why it remains challenging even with
positive-class weighting. SPADE~\cite{spade} provides a segment-level relation
formulation, while BROS~\cite{bros} uses sparse linking between entity-initial
positions. Both provide alternatives to dense token-to-token relation scoring.

\subsection{Comparison to Baselines}

The final corrected V4 reaches Regular F1 of $0.392$ for text, $0.446$
for location, and $0.345$ for text+location. It combines the corrected data and
supervision pipeline with correct scheduler accounting, full-state resumption,
class-aware entity evaluation, validation pair-F1 checkpoint selection, and
aligned early stopping. Because these corrections are applied together, their
individual effects are not isolated.

The reconstructed Mistral baseline remains stronger in the official comparable
Regular modes, with text F1 of $0.598$ and text+location F1 of $0.521$. V4 and
the reconstructed Mistral baseline were evaluated on the same 581 locally
available pages with the same evaluator, so these internal results are directly
comparable. The published Regular values of $0.659$ and $0.611$ are included
only as contextual information. The generative and discriminative systems have different strengths: Mistral
achieves higher Regular pair scores in the comparable benchmark modes, while
the discriminative system cannot emit text that is absent from the extracted
input. V4 provides the final discriminative result of this study but remains
below the reconstructed Mistral baseline on complete Regular pair extraction.

The direct V4 decoder predicts Regular pairs only. Its direct All
text+location F1 is $0.229$. The unchanged recovery method converts unmatched
VALUE spans into unkeyed outputs and unmatched KEY spans into unvalued outputs.
This raises All text+location F1 to $0.261$, with $0.162$ for unkeyed and
$0.276$ for unvalued items. Recovery improves special-category coverage but
does not change any Regular prediction or the Regular F1 of $0.345$.

\subsection{Per-Cluster Performance}

The frozen Stage~2 assignment divides the prepared test subset into 319 pages
in Cluster~0, 213 pages in Cluster~1, and 49 pages without usable annotation
geometry. Of these, 304 and 179 pages, respectively, have scorable Regular
ground truth. The geometry-unavailable pages do not have scorable Regular
ground truth and are excluded from the category macro average.

\begin{table}[ht]
  \centering
  \caption{Regular text+location macro F1 by page-level annotation-geometry
  cluster. The clusters are geometry regimes, not document types.}
  \label{tab:stage4-cluster}
  \begin{tabular}{@{}lcc@{}}
    \toprule
    \textbf{Model} & \textbf{Cluster 0} & \textbf{Cluster 1} \\
    \midrule
    Reconstructed Mistral & 0.569 & 0.429 \\
    V4                    & 0.305 & 0.410 \\
    \bottomrule
  \end{tabular}
\end{table}

V4 performs better in Cluster~1 than in Cluster~0 in text, location, and joint
evaluation. Mistral has the opposite ordering. Annotation geometry is
associated with system performance. The observed association differs between
the two systems: V4 and Mistral show opposite performance ordering across the
two geometry clusters. It is not evidence of a
causal effect, and it does not show performance differences between semantic
document types. The modest training silhouette score also limits strong claims
about cluster separation.

\section{Root-Cause Analysis}
\label{sec:root-cause-analysis}

The progression shows that architectural changes are difficult to assess when
the supervision pipeline contains defects. The diagnostics identified three
interacting pipeline failures affecting span construction and relation
supervision.  Because V4 introduced the corrections jointly, their separate
numerical contributions were not measured; the following mechanisms are
described without ranking individual effects.

\paragraph{Missing teacher forcing.}
The original predicted-span variant formed spans from \texttt{argmax} entity
predictions while relation labels came from ground-truth annotations. This
created a mismatch between the spans supplied to the linker and the
ground-truth relation supervision. V4 uses ground-truth spans during training
and predicted spans during inference. This introduces the usual
train--inference difference associated with teacher forcing.

\paragraph{Bounding-box filter.}
The strict $(x_1 < x_2) \wedge (y_1 < y_2)$ predicate could reject valid tokens
whose quantised LayoutLMv3 boxes had zero width or zero height. On 200 pages
from the seed-42 validation split, it removed approximately 47.5\% of predicted
KEY tokens and 48.5\% of predicted VALUE tokens. Only 81 of 200 pages retained
both strict key and value spans.

\paragraph{Cross-product link labels.}
The label generator created all-pairs links between key and value words rather
than a single representative link per KVP. This made the span-level relation
supervision too dense to represent one relation per annotation.

\noindent The cross-product issue became easier to diagnose after span-level
collapse, because the resulting span-level target matrix was expected to
represent sparse annotated relations. The span-level reformulation therefore
also served as a useful diagnostic change in addition to reducing the candidate
space.

\section{Error Analysis}
\label{sec:error_analysis}

\subsection{V4 Qualitative Errors}

The selected examples in Table~\ref{tab:v4-qualitative-errors} illustrate
several error types. They are consistent with the gap between the entity and
pair metrics: a model can identify plausible entities but still select the
wrong relation or miss a relation. The cases are not a statistical sample and
do not prove their frequency. Some source annotations are also semantically
ambiguous, so individual cases require cautious interpretation.

\subsection{Validation Diagnostics for V1 and V2}

The fixed seed-42 validation subset is used to inspect link scores, predicted
entity spans, bounding-box filtering, and relation-label construction. The
token-level V1 formulation has a dense candidate space. Span-level V2 reduces
that space, but predicted-span alignment, the strict bounding-box filter, and
cross-product labels prevent a clean assessment of linking granularity. These
checks are internal diagnostics and are not official page-macro benchmark
results.

For the preserved V2 checkpoint, 1{,}481 candidate-pair logits have a median
of $-6.93$ and a mean of $-7.59$. Only 36 candidate pairs and 20 best-per-key
predictions have sigmoid probability at least $0.5$.

The corrected label-pipeline audit on the first 20 validation pages contains
79 regular ground-truth pairs, 53 recovered representative word-level links,
453 positive token-pair cells after subword expansion, and 44 positive
span-pair labels after span collapse.

\section{Relation-Modeling Context and Experimental Scope}
\label{sec:relation_context}

The experimental comparisons in this thesis remain anchored to the KVP10k
Mistral baseline and the LayoutLMv3 encoder with a projected scaled-dot-product linker with spatial features. As discussed
in Section~\ref{sec:experimental_scope}, other relation-extraction methods were
available before the project started, but implementing them would have required
separate model designs and full experimental pipelines. SPADE~\cite{spade},
BROS~\cite{bros}, GeoLayoutLM~\cite{geolayoutlm},
KVPFormer~\cite{kvpformer}, and LayoutPointer~\cite{layoutpointer} are therefore
included as methodological context and future alternatives rather than as
additional experimental baselines.

SPADE and BROS are evaluated on document-understanding benchmarks other than
KVP10k, so their reported numbers are not directly comparable with the results
in this thesis. SPADE's segment-level dependency formulation and BROS's sparse
entity-initial linking objective nevertheless provide relevant methodological
context for reducing dense token-to-token relation scoring.

\section{Limitations and Future Work}
\label{sec:limitations}

\subsection{Current Limitations}

Several limitations constrain the interpretation of the results. First, the
prepared development pool contains 5{,}389 usable pages from the 9{,}656 unique
pages in the official KVP10k training split, or 55.8\%. The fixed seed-42 split
then uses 4{,}851 pages for model training and 538 for validation. Some PDFs
were unavailable or broken, and others had no usable text layer. The pipeline uses PyMuPDF native PDF text extraction rather
than OCR, so scanned pages without a text layer are excluded.

Second, the reconstructed Mistral baseline is a re-implementation rather than
an exact reproduction. Its LoRA rank, prompt details, input extraction, and
training coverage are not confirmed to match the published system. Third, V4
uses constant blank visual input. It evaluates textual and layout features but
not document-specific image features. Fourth, inputs are limited to 512 tokens
and one page. Truncation can remove later content, and the current architecture
would not model cross-page relations when applied to multi-page documents. Finally, the qualitative cases are
illustrative, and annotation ambiguity can affect their interpretation.

\subsection{Future Directions}

A primary direction for future work is improved relation decoding. Structured
one-to-one or constrained matching could reduce incompatible local link
decisions. Hard-negative
training could focus on plausible neighbouring values, as illustrated by the
selected V4 examples. Explicit row and column features can add document
structure that is not captured by pairwise displacement alone. These changes
can also be compared with segment-graph and relation-focused methods such as
SPADE~\cite{spade}, BROS~\cite{bros}, GeoLayoutLM~\cite{geolayoutlm},
KVPFormer~\cite{kvpformer}, and LayoutPointer~\cite{layoutpointer}.

The output design also requires improvement. Learned unkeyed and unvalued
classes could replace heuristic recovery, while confidence calibration could
improve the reliability of the fixed decision threshold. Better segmentation and longer
input handling are needed for short values, long pages, and multi-page
documents. Finally, a controlled experiment with rendered document images is
needed to measure the contribution of genuine visual features. All these items
are future work; they were not implemented in V4.
